\chapter{Arthroscopy}

\section*{Overview}
Arthroscopy is a minimally invasive surgical procedure on a joint in which an examination and sometimes treatment of damage is performed using an arthroscope. The exact approach, setting, and preparation depend on the indication, anatomy, and the patient's health.

\section*{Alternative Names}
Joint scope, Keyhole joint surgery

\section*{Principle}
A narrow fiberoptic scope with a camera is inserted into a joint through a small incision and the joint is distended with fluid. The interior is viewed on a monitor while instruments passed through additional small ports diagnose and treat damage, minimizing trauma to surrounding structures.
\section*{History}
Kenji Takagi performed the first arthroscopy on a cadaver knee in 1918. Masaki Watanabe developed modern arthroscopic instruments and techniques in the 1950s.

\section*{Why It Is Done}
Ordered to diagnose and treat joint pain, swelling, or instability. Commonly performed on knees, shoulders, elbows, ankles, hips, and wrists.

\section*{Is This a Primary or Secondary Test or Procedure}
Primary surgical standard for joint-preserving procedures (e.g., meniscus repair, ACL reconstruction).

\section*{How It Is Performed}
Performed under general, regional, or local anesthesia. Small incisions (portals) are made around the joint. The scope and surgical instruments are inserted to inspect and repair the joint.

\section*{Preparation}
If sedation or anesthesia is planned, follow the facility's fasting instructions. Review medicines, allergies, and health conditions with the surgical team; do not stop anticoagulants or other medicines unless given a specific plan. Arrange transportation and help at home when mobility may be limited.

\section*{Contraindications and Precautions}
Active skin infection over the joint, or severe joint stiffness (ankylosis) preventing scope maneuverability.

\section*{Risks}
Joint infection (septic arthritis), hemarthrosis (bleeding in the joint), nerve damage, or blood clots (DVT).

\section*{Aftercare and Recovery}
Keep incisions clean and dry. Apply ice and elevate the joint to reduce swelling. Follow weight-bearing instructions.

\section*{Outcome and Follow-up}
Visualizes cartilage tears, ligament tears, loose bodies, or synovitis. Surgical repair is performed during the same procedure.

\section*{Expected Results}
Intact articular cartilage; normal ligaments and meniscus; no loose bodies or inflammation.

\section*{Follow-up}
Physical therapy, wound check in 7-10 days, and gradual return to activity.

\section*{When to Seek Medical Care}
Follow the procedure team's specific instructions. Seek urgent medical assessment for severe or worsening pain, uncontrolled bleeding, fever or signs of infection, trouble breathing, chest pain, fainting, new weakness or confusion, inability to urinate, or any rapidly worsening symptom. The appropriate warning signs depend on the procedure and the patient's medical condition.

\section*{References}
\begin{enumerate}
  \item MedlinePlus. Arthroscopy. U.S. National Library of Medicine; 2025.
  \item Current Medical Diagnosis and Treatment. McGraw-Hill; 2025.
\end{enumerate}

\clearpage
